\documentclass[11pt]{article}
\usepackage[a4paper,margin=1in]{geometry}
\usepackage{amsmath,amssymb,amsthm,mathtools}
\usepackage{siunitx}
\usepackage{graphicx}
\usepackage{xcolor}
\usepackage{microtype}
\usepackage{booktabs}
\usepackage{hyperref}
\usepackage[nameinlink,capitalise]{cleveref}

\hypersetup{
  colorlinks=true,
  linkcolor=blue!50!black,
  urlcolor=blue!60!black,
  citecolor=blue!60!black
}

\newcommand{\R}{\mathbb{R}}
\newcommand{\C}{\mathbb{C}}
\newcommand{\Z}{\mathbb{Z}}
\newcommand{\e}{\mathrm{e}}
\newcommand{\im}{\mathrm{i}}
\newcommand{\T}{\mathsf{T}}
\newcommand{\dd}{\,\mathrm{d}}
\newcommand{\half}{\tfrac{1}{2}}
\newcommand{\elll}{\ell}  % to avoid confusion with \ell in math fonts
\newcommand{\figpath}{.}  % change to ./figs or similar if needed

\title{\vspace{-2ex}\Large Temporal Envelopes, Decimated AR(2) Dynamics, and Physics-Native Normalization:\\
From Floquet Quantum Mechanics to ATLAS \texorpdfstring{$\varphi$}{phi}-Histograms}
\author{}
\date{\today}

\begin{document}
\maketitle

\begin{abstract}
We develop and test a discrete-time, discrete-space framework in which coarse observables evolve by
a second-order autoregressive (AR(2)) law across a rank lattice. We show three routes by which such
dynamics emerge from standard quantum mechanics (QM): (i) eliminating an internal two-component
Floquet ``coin'' in a discrete-time quantum walk, (ii) two-frequency dominance of expectation-value
time series, and (iii) coarse-grained Mori--Zwanzig closures. We then show that, for any one-dimensional
second-order medium, \emph{decimating} by stride $\ell$ induces a universal tri-site recurrence
$x_{k+2\ell} - V_\ell x_{k+\ell} + Q^\ell x_k = 0$ where $V_\ell=\mathrm{tr}(M^\ell)$ and $Q=\det M$
are the trace and determinant of the $\ell$-step monodromy $M^\ell$. The coefficients $V_\ell$ obey a
Lucas recursion from the one-step pair $(V,Q)$.

On top of this identity we build a physics-native normalization/decoding calculus on each residue
track $k=r+\ell t$ using a local carry/borrow rule that provably terminates and is confluent. We
define a Chebyshev--Lucas checksum $I(d)=\sum_t d_t r^t$ (for a root $r$ of $z^2-V_\ell z + Q^\ell=0$)
that is invariant under the local moves. We implement the full pipeline on ATLAS public \emph{PHYSLITE}
$\varphi$-histogram data: eventization, per-track normalization, modular checksum validation, and
stride-$\ell$ recomposition. In a representative configuration (stride $\ell=16$ on electron-like
histograms with $\varepsilon=10$ event quantum) the RMS violation of the stride-$\ell$ identity
drops from $\sim\!\num{1.50e7}$ to $\sim\!\num{2.83e6}$ after normalization (factor $\approx5$).
Modular checksum tests (two 61--64-bit primes) pass bit-for-bit.

We also include a superdeterministic envelope module for Bell--CHSH simulations in which a hidden
seam phase $\lambda$ biases setting choices with a tunable measurement-dependence parameter $\eta$,
with an optional $\pi$-symmetrization enforcing no-signaling. This allows principled tests of how
much alignment between envelope phases and setting RNGs would be required to reproduce quantum-like
violations in a local model.
\end{abstract}

\section{Motivation and Overview}

We ask whether a minimal spatiotemporal law,
\begin{equation}
\label{eq:ar2-full}
x_{m+2}(j) = V\,x_{m+1}(j) - Q\,x_m(j) + r\!\left(\e^{\im\xi} x_{m+1}(j{+}1) + \e^{-\im\xi} x_{m+1}(j{-}1)\right),
\end{equation}
can capture the long-wavelength, coarse-grained kinematics of observables built from complex
microscopic dynamics. Here $m\in\Z$ is a stroboscopic time index and $j\in\Z_L$ labels a rank/angle
position (e.g., $\varphi$-bin). For plane waves $x_m(j)=\e^{\im(kj-\omega m)}$, \eqref{eq:ar2-full}
implies
\begin{equation}
\label{eq:dispersion}
2\cos\omega \;=\; V + 2r\cos(k+\xi)\qquad(Q=1).
\end{equation}
Boundary twists (APBC, M\"obius-like holonomy) shift $k\mapsto k+\pi/L$ and open a spectral gap.
Within this law we can \emph{measure} $(V,Q,r,\xi)$ from data via $k$--$\omega$ overlays.

A central structural claim is that---independently of the microscopic origin---\emph{decimating} a
one-dimensional second-order chain by stride $\ell$ always yields a tri-site identity
\begin{equation}
\label{eq:decimated}
x_{k+2\ell} - V_\ell\,x_{k+\ell} + Q^\ell x_k = 0,\qquad
V_\ell=\mathrm{tr}(M^\ell),\ \ Q=\det M,
\end{equation}
where $M$ is the one-step transfer (monodromy) matrix. The $V_\ell$ obey the Lucas/Dickson recursion
$V_0=2,\ V_1=V,\ V_{n+1}=V V_n - Q V_{n-1}$. This identity underlies both our normalization calculus
and a strong multi-resolution self-consistency check for fitted parameters.

\paragraph{Contributions.} (i) Three independent derivations of the AR(2) law from textbook QM;
(ii) a decimation theory that links one-step and stride-$\ell$ recurrences; (iii) a per-track
normalization/decoding algorithm (deterministic carries) with a physics-native checksum invariant;
(iv) a practical pipeline tested on ATLAS \emph{PHYSLITE} $\varphi$-histograms; (v) a superdeterministic
envelope module for controlled measurement dependence in CHSH tests.

\section{AR(2) from Standard Quantum Mechanics}
\label{sec:qm-to-ar2}

\paragraph{Floquet/coin route (exact).}
For a discrete-time walk with two-component ``coin'' $\Phi_{m+1}(k)=U(k)\Phi_m(k)$ and unitary
$U(k)\in\mathrm{U}(2)$, the characteristic polynomial of $U$ is
$\lambda^2 - \tau(k)\lambda + \det U(k)=0$. Eliminating one component yields
\begin{equation}
z_{m+2}(k)-\tau(k)\,z_{m+1}(k)+\det U(k)\,z_m(k)=0.
\end{equation}
After a global phase gauge, $\det U(k)=1$ and $\tau(k)=V+2r\cos(k+\xi)$, giving exactly
\eqref{eq:ar2-full} with $Q=1$ on inverse Fourier transform.

\paragraph{Two-frequency dominance (exact for two lines).}
If an observable has two Bohr lines $\Omega_{1,2}$ in the sampled sequence
$x_m=\langle A(m\Delta t)\rangle$, then $x_m$ obeys
$x_{m+2}- (e^{-\im\Omega_1\Delta t}+e^{-\im\Omega_2\Delta t})x_{m+1} + e^{-\im(\Omega_1+\Omega_2)\Delta t}x_m=0$,
an AR(2) with complex $(V,Q)$. With more lines the AR(2) is the best linear predictor on windows.

\paragraph{Mori--Zwanzig (approximate with innovations).}
Projecting unitary dynamics onto a slow subspace yields a second-order equation with memory. With
a short-memory (two-pole) closure and stroboscopic sampling, one obtains \eqref{eq:ar2-full} with
an innovation~$\eta_m$ that our high-pass whitening largely removes.

\paragraph{Continuum limit (Klein--Gordon).}
Setting $\omega\Delta t\ll1$, $ka\ll1$ and expanding \eqref{eq:dispersion} gives
$\omega^2 \approx \mu^2 + c^2k^2$ with parameters determined by $(V,r,\xi)$ and lattice spacings;
this is the standard relativistic dispersion in the long-wavelength limit.

\section{Decimation and Chebyshev--Lucas Structure}
\label{sec:decimation}

Consider a one-dimensional second-order medium with one-step monodromy $M\in\C^{2\times2}$ and
$(V,Q)=(\mathrm{tr}M,\det M)$. By Cayley--Hamilton, $M^2 - V M + Q I = 0$. Sampling every $\ell$
sites we work with $M^\ell$ and obtain the \emph{decimated} tri-site identity
\eqref{eq:decimated}. The coefficient $V_\ell=\mathrm{tr}(M^\ell)$ satisfies the Lucas recursion
$V_0=2$, $V_1=V$, $V_{n+1}=V V_n - Q V_{n-1}$.

\paragraph{Checksum.} Let $r$ be a root of $z^2 - V_\ell z + Q^\ell=0$. For a (finite-support)
digit string $d=\{d_t\}$ on a residue track $k=r_0+\ell t$, the quantity
\begin{equation}
I(d)=\sum_t d_t\,r^t
\end{equation}
is invariant under the local tri-site update $(d_{t-1},d_t,d_{t+1})\mapsto(d_{t-1}+Q^\ell q,\ d_t-V_\ell q,\ d_{t+1}+q)$.
We use this as a physics-native parity check and as the basis of a small-window trellis decoder.

\section{Complex-Base Normalization on Residue Tracks}
\label{sec:normalization}

Let the complex base be $B=b\,\e^{\im2\pi/\ell}$ where $b$ solves $b^2-Pb+Q=0$. Grouping indices by
residue $r\in\{0,\dots,\ell-1\}$ gives tracks $k=r+\ell t$. For even $\ell$ with $Q^\ell>0$ the
canonical digit alphabet is nonnegative $\{0,\dots,V_\ell-1\}$; for odd $\ell$ with $Q^\ell<0$ we
use a balanced alphabet $\{-a,\dots,+a\}$, $a=(V_\ell-1)/2$. A rightmost-first rule with the local
tri-site move terminates under a decreasing potential $E_q=\sum_t d_t q^t$ for $q\in(1,b^\ell)$,
leading to a unique normal form per track. In practice we normalize each track independently and
verify $I(d)$ modulo large primes to avoid floating overflow.

\section{Superdeterministic Envelope Module}
\label{sec:sd}

We model a hidden seam phase $\lambda\in[0,2\pi)$ that biases CHSH settings with strength $\eta\in[0,1]$:
the preferred setting at a site is the one maximizing $|\cos(\lambda-\theta)|$, chosen with probability
$(1+\eta)/2$ (the other with $(1-\eta)/2$). Deterministic local outcomes are
$A=\mathrm{sgn}[\cos(\lambda-a)]$, $B=-\mathrm{sgn}[\cos(\lambda-b)]$. To enforce no-signaling we
$\pi$-symmetrize $\lambda$ after settings selection (replace $\lambda\to\lambda$ or $\lambda{+}\pi$
with equal probability), which flips both marginals while preserving the product $AB$. We report
CHSH $S$, an empirical measurement-dependence index $\hat M$ (TV distance between $p(\lambda|a,b)$
histograms), and marginal drift diagnostics.

\section{Data and Pipeline}
\label{sec:data}

We use ATLAS public \emph{PHYSLITE} $\varphi$-histogram NPZ files of shape $(T,L)$ (time windows by
$\varphi$-bins). The pipeline steps are:

\begin{enumerate}
\item \textbf{Eventization:} For each rank $j$, integrate time by $E_j=\sum_m\max(Y_{m,j},0)$
or $E_j=\sum_m|Y_{m,j}|$ and quantize with an event quantum $\varepsilon$ to $N_j=\lfloor E_j/\varepsilon\rfloor$.
\item \textbf{Track split:} Split $j$ into residue classes $j=r+\ell t$.
\item \textbf{Normalization:} For each track, apply the local tri-site carry/borrow until all digits
lie in the legal alphabet; count unit carries.
\item \textbf{Checksum:} Verify modular equality of $\sum_j N_j r^j$ and $\sum_{r,t} d^{(r)}_t r^{r+\ell t}$
for two large primes (e.g., $2^{61}{-}1$ and $2^{64}{-}59$).
\item \textbf{Stride-$\ell$ recomposition:} Reassemble a one-dimensional lattice $x[j]$ and check
the residuals of $x_{j+2\ell}-V_\ell x_{j+\ell}+Q^\ell x_j$.
\end{enumerate}

\paragraph{Commands.} Example invocation of the bridge:
\begin{verbatim}
python complex_base_bridge.py \
  --npz PHYSLITE_phiHist_hp256_ell64_Y_ell64.npz \
  --ell 16 \
  --P 1 --Q -1 \
  --mode abs_cumsum --eps 10 \
  --alphabet auto \
  --out out_eps10_ell16
\end{verbatim}
Recomposition/AR(2) check:
\begin{verbatim}
python recompose_spatial_check.py \
  --digits-raw out_eps10_ell16/digits_raw.csv \
  --digits-norm out_eps10_ell16/digits_normalized.csv \
  --ell 16 --P 1 --Q -1 --out spatial_eps10_ell16
\end{verbatim}

\section{Results}
\label{sec:results}

\paragraph{Checksum validation.}
Modular parity checks matched exactly for two independent large primes, providing a strong pass
without floating overflow. Per-track invariants $I(d)$ are preserved to machine precision.

\paragraph{Stride-$\ell$ residual reduction.}
For $\ell=16$ (Lucas $V_{16}=2207$, $Q^{16}=1$), eventization with $\varepsilon=10$ produced
nontrivial carries on most tracks. Recomposition residual RMS for the stride-$\ell$ identity decreased
from $\sim\!\num{1.50e7}$ (raw digits) to $\sim\!\num{2.83e6}$ (normalized digits), a factor $\approx5$ improvement.

\paragraph{Visualization.}
\cref{fig:komega-demo} shows a $k$--$\omega$ overlay of a synthetic transport stack built from the
normalized lattice to illustrate how ridges would appear with time dynamics. \cref{fig:residuals}
plots the stride-$\ell$ residuals before and after normalization.

\begin{figure}[t]
  \centering
  \includegraphics[width=0.9\linewidth]{\figpath/komega_demo.png}
  \caption{$k$--$\omega$ overlay (synthetic transport demo) built from normalized digits.}
  \label{fig:komega-demo}
\end{figure}

\begin{figure}[t]
  \centering
  \includegraphics[width=\linewidth]{\figpath/residuals.png}
  \caption{Stride-$\ell$ residuals ($\ell=16$): raw (blue) vs.\ normalized (orange).}
  \label{fig:residuals}
\end{figure}

\section{Discussion and Outlook}
\label{sec:discussion}

\paragraph{Internal consistency.}
The Chebyshev--Lucas structure provides a stringent self-consistency ladder across strides.
Fitting a one-step pair $(V,Q)$ from $k$--$\omega$ overlays, we must recover $V_\ell$ by the
Lucas recursion and $Q^\ell$ by exponentiation. Disagreement flags model or preprocessing issues.

\paragraph{Physical interpretation and limits.}
Our derivations place the AR(2) law inside orthodox QM via two-component Floquet elimination or
two-line dominance; decimation is a structural identity of second-order chains. What remains
\emph{speculative} is whether these coarse rails and their topological locks provide a faithful
parametrization of high-energy dynamics. The present results show a mathematically precise way to
regularize and validate rank-lattice data, not a claim of fundamental ontology.

\paragraph{Falsifiable checks.}
(i) Lucas consistency across strides for $(V,Q)$ from real overlays; (ii) controlled opening of
holonomy gaps (APBC) in engineered Floquet systems; (iii) envelope-driven CHSH tests where the
measured $S$ is mapped against the empirical measurement-dependence index $\hat M$.

\paragraph{Near-term experimental routes.}
Photonic split-step walks and cold-atom Floquet realizations allow direct programming of $U(k)$,
phase defects (holonomy), and time strobing. Our pipeline can be applied out of the box to their
time-resolved intensity maps.

\section*{Reproducibility: Minimal Code Listings}

\paragraph{Stride-$\ell$ coefficients via Lucas.}
\begin{verbatim}
def lucas_Vn(P, Q, n):
    if n == 0: return 2.0
    if n == 1: return float(P)
    Vm1, V = 2.0, float(P)
    for _ in range(2, n+1):
        Vm1, V = V, P*V - Q*Vm1
    return V

def stride_coeffs(V1, Q1, ell):
    return lucas_Vn(V1, Q1, ell), (Q1**ell)
\end{verbatim}

\paragraph{Per-track normalization (deterministic).}
\begin{verbatim}
def normalize_track(d, Vell, Qell, balanced=False):
    d = {int(k): int(v) for k,v in d.items() if v!=0}
    V = int(round(Vell))
    carries = 0
    a = (V-1)//2
    def inD(x): return (-a <= x <= a) if balanced else (0 <= x <= V-1)
    while True:
        bad = [t for t,v in d.items() if not inD(v)]
        if not bad: break
        t = max(bad); v = d[t]
        if balanced:
            q = int(np.floor((v + a)/V)); d[t] = v - q*V
        else:
            if v >= V: q = v // V; d[t] = v - q*V
            else:      q = (-v + V - 1)//V; d[t] = v + q*V
        d[t-1] = d.get(t-1,0) + int(Qell)*q
        d[t+1] = d.get(t+1,0) + q
        carries += abs(q)
        # prune zeros ...
    return d, carries
\end{verbatim}

\paragraph{Modular checksum.}
\begin{verbatim}
def checksum_global_modp(N_rank, tracks_norm, ell, r, primes):
    out=[]; r=int(r)
    for p in primes:
        rp = r % p; r_ell = pow(rp, ell, p)
        Iev=0; Idg=0
        for j, Nj in enumerate(N_rank):
            if Nj: Iev = (Iev + (Nj % p) * pow(rp, j, p)) % p
        for r0, d in tracks_norm.items():
            rr0 = pow(rp, r0, p)
            for t, val in d.items():
                if val: Idg = (Idg + (val%p) * rr0 * pow(r_ell, t, p)) % p
        out.append(dict(p=p, equal=(Iev==Idg)))
    return out
\end{verbatim}

\section*{Data Availability and Acknowledgements}
The NPZ files used here are public ATLAS \emph{PHYSLITE}-derived $\varphi$-histograms prepared by
the authors from the Open Data portal. Intermediate CSVs and figures referenced in this manuscript
are produced by the scripts distributed alongside this \LaTeX\ source. This document contains no
proprietary or personally identifiable information.

\end{document}
